\chapter{Consensus and Replication}

CharlotteOS exposes distributed consensus as a capability service: a replicated
state machine reached through the same endpoint IPC used by local services.

The test suite covers single-voter term recovery from NVMe and cross-machine
catalog replication and invocation through \texttt{dns}. Each operational node
runs exactly one durable Raft member inside \texttt{dns}; membership, names,
deployments, keys, and cluster events share that member's ordered log. The
multi-guest admission test observes it through discovery and exercises
\texttt{JOIN}, joint consensus, and \texttt{FINALIZE}. The separate
\texttt{raft.elf} process is network-disabled and retained only by the NVMe
term-recovery fixture. Power-loss behavior within a live quorum is not yet
validated end to end.

The core follows the shared Graft model for
normalized peer identity, voters and learners, joint-consensus quorums,
replicated \texttt{JOIN}/\texttt{JOINT}/\texttt{FINALIZE} commands, leader
no-op entries, current-term commit rules, linearizable-read contact quorums,
and membership-bearing snapshots. These paths have focused regression tests. The
\texttt{dns} service exposes both client commands and membership
administration end to end through one concrete state machine and log.

\section{Raft as a capability service}

The Raft consensus protocol provides a replicated, fault-tolerant state machine
across a cluster of nodes. In CharlotteOS, the operational Raft node is owned
by the per-node DNS EL0 process, which:

\begin{itemize}
    \item Registers the DNS endpoint and a \texttt{"raft-<id>"}
        administration/status endpoint backed by the same in-process node.
    \item Learns transient peer MAC routes through discovery; voter authority
        changes only through committed membership commands.
    \item Carries VoteRequest, AppendEntries, InstallSnapshot, and admission
        messages through the reliable-message service.
    \item Uses completion-queue waits for time --- a bounded
        \texttt{CQ\_WAIT} drives election and heartbeat checks without a
        detached timer completion or a busy polling loop.
    \item Applies membership commands and DNS catalog commands to the same
        durable log and state-machine snapshot.
\end{itemize}

The Raft core (\texttt{catten-graft}) is a generic, transport-agnostic library.
It implements the \texttt{RaftTransport} trait, which abstracts peer
communication behind a few methods:

\begin{verbatim}
trait RaftTransport {
    fn send_vote_request(&self, request: VoteRequest);
    fn send_append_entries(&self, request: AppendEntries);
    fn send_install_snapshot(&self, request: InstallSnapshot);
    fn poll_completions(&self) -> Vec<RpcCompletion>;
}
\end{verbatim}

The operational CharlotteOS implementation uses the reliable message layer
between DNS replicas; the consensus core does not change.

\section{Intended properties relative to socket-based Raft}

\subsection{Authority, not addresses}

Once the cluster is running, peers are identified by name service entries, not
\texttt{IP:port}. A node that has joined the cluster receives connection caps
for its peers through the name service:

\begin{verbatim}
let peer_1: ConnectionCap = name_service.lookup("raft-node-1")?;
let peer_2: ConnectionCap = name_service.lookup("raft-node-2")?;
\end{verbatim}

No IP configuration. No DNS. The authority to communicate with a peer is the 
capability; without it, you cannot even address the peer.

\subsection{Transport transparency}

The transport is behind the \texttt{RaftTransport} trait. A connection capability
routes through local IPC (same machine), Ethernet frames (same network), or a
future RDMA transport (high-performance cluster) --- the consensus core never
changes. This means:

\begin{itemize}
    \item Development and testing can run a multi-node Raft cluster on one
        machine with local IPC.
    \item Deployment can switch to cross-machine networking by changing how
        peers are resolved, not by changing the Raft code.
    \item The cluster can be heterogeneous: some peers communicate via shared
        memory, others via Ethernet.
\end{itemize}

\subsection{Zero-copy log replication}

Locally, log payloads can transfer via Move memory objects: the kernel flips page-table
ownership instead of copying bytes. A follower that receives a Move attachment
gets the exact same physical pages the leader held --- no intermediate buffer,
no copy. 

The current local transport uses one moved page per RPC and batches the largest
AppendEntries prefix whose protobuf encoding fits. This is a 4\,KiB transport
attachment limit, not a persistent-object or Raft-log limit. A network
transport can use different framing without changing the protobuf contract.

\subsection{Capability-based security}

A client invokes the Raft service only if it possesses a connection cap. The
service manager attenuates rights: a typical client receives only
\texttt{CALL} for \texttt{OP\_CLIENT\_COMMAND}, not for
\texttt{OP\_ADD\_SERVER} or \texttt{OP\_REMOVE\_SERVER}. There is no ambient
authority, no IP-based ACL, and no ``any process on the cluster network can
propose commands.''

\subsection{Failure isolation}

Each DNS/Raft replica runs in a separate protection domain. A crash in one node
does not affect the others. The isolated storage fixture demonstrates recovery
of term/vote state across explicit process teardown and restart; restarting a
DNS replica and recovering its catalog inside a live quorum remains future
integration work. The generic service lifecycle separately demonstrates that
stale connections fail and clients can re-lookup a replacement generation.

\section{Persistence and launch configuration}

The operational DNS node requires persistent storage. Stable object identifiers
derived from cluster and node identity hold term/vote state, log entries, and
snapshot data; mutations are followed by an NVMe flush.

The cluster mnemonic and election timeout remain in the launch manifest. A
fresh identity bootstraps as a durable singleton. Discovery supplies MAC routes,
the lexicographically larger singleton applies to the smaller anchor, and
active membership changes only through the Raft log and snapshot envelope.
The distinct storage fixture verifies one generic node across process restart;
the two-guest DNS test verifies the operational cluster.

\section{Cluster bootstrap --- the seed problem}

A new node must contact at least one existing cluster member. CharlotteOS
expresses the seed as a service name or capability rather than an IP address.

\subsection{Ethernet discovery and deterministic admission (current)}

Nodes broadcast their durable identity on the local Ethernet segment. DNS
records every discovered identity-to-MAC route in its relmsg transport. Among
fresh singletons, only the lexicographically larger identity applies to the
smaller anchor. Before sending the request it durably enters a non-voting
joining posture; the anchor then commits \texttt{JOIN}, catches the node up,
commits the joint configuration, and finalizes it.

\subsection{Pre-shared cluster capability (small kernel change)}

A cluster administrator provisions a ``cluster capability'' --- a connection
to a known cluster member --- and delivers it to new nodes via the config page.
No name lookup required; the node already holds a first-class connection.

\subsection{Ethernet broadcast discovery (zero-configuration)}

On a single Ethernet segment, a joining node broadcasts a frame with a
well-known EtherType carrying a ``Raft peer discovery'' payload. Existing
members respond with their service names. The joining node looks them up
through the name service to obtain attested connection caps. This works
without any seed configuration.

\subsection{Distributed name service (implemented)}

The \texttt{dns} service realizes this end state: its \texttt{CatalogMachine}
is a \texttt{StateMachine} implementation on the Raft substrate, and a new
node needs only a single seed to discover all peers through the consistent
name registry. There remains one process-facing name-service instance per
machine; those instances apply the same committed registry state. The
\texttt{ns} node-local registry and the \texttt{dns} cluster catalog remain
distinct services.

\section{The generic StateMachine trait}

Raft replicates an opaque command log. The meaning of those commands is provided
by a pluggable \texttt{StateMachine} trait:

\begin{verbatim}
trait StateMachine {
    fn apply(&mut self, command: Vec<u8>) -> Result<Vec<u8>, ApplyError>;
    fn snapshot(&self) -> Vec<u8>;
    fn restore(&mut self, snapshot: Vec<u8>);
}
\end{verbatim}

Different services are different state machines on the same Raft substrate:

\begin{itemize}
    \item Distributed name service --- commands are
        ``register(name, owner, route, generation)'' and
        ``unregister(name, owner, generation)''. Unregistration is rejected
        unless both owner and generation still match the active entry.
        The owning DNS process watches the local endpoint through a waitable
        kernel completion. Endpoint death proposes this same fenced tombstone;
        a follower forwards an authenticated, idempotent request to its known
        leader and retries across leadership changes. A stale death notification
        therefore cannot remove a replacement generation.
        Scalar registration adopts an existing node-local service connection
        through a non-blocking local lookup. This lets DNS route to and observe
        the endpoint without granting ordinary lookup clients minting rights;
        names without a local endpoint remain catalog-only registrations.
        The state maps names to owning-node identities, routable service
        identities, generations, and policy metadata. Replicated, it makes
        service discovery available cluster-wide; this is the implemented
        \texttt{dns} service.
    \item Distributed lock service --- commands are ``acquire(lock\_id)''
        and ``release(lock\_id)''. The state is a set of held locks.
    \item Cluster configuration store --- commands are arbitrary
        key-value writes. The state is a key-value map.
\end{itemize}

None of these state machines know about sockets, IP addresses, or TCP. They
implement \texttt{apply()} and \texttt{snapshot()}. The Raft substrate handles
replication, leader election, and fault tolerance.

Raft RPCs and DNS control messages use one serialized relmsg queue per peer.
Because QEMU may acknowledge that queue slowly, consecutive queued
AppendEntries requests and their serially ordered responses are coalesced, and
heartbeats are generated only at a bounded fraction of the election timeout.
Non-Raft control messages retain their relative FIFO order but pass queued,
supersedable append traffic. A two-second stop-and-wait retry lease is shorter
than the DNS remote-call deadline, preventing one lost heartbeat from holding
the peer's only transmission slot long enough to starve service retraction,
placement, or invocation traffic.

The DNS process also never waits synchronously for an arbitrary local service
while running this reactor. Name lookup and local invocation are retained as
bounded pending operations and polled on later turns. Consequently a failed
service cannot stop Raft heartbeats or relmsg receive draining. A timeout before
execution is reported as a retryable lookup failure; after invocation has been
submitted it is reported as an uncertain outcome. Duplicate remote calls are
suppressed both while pending and in the acknowledged-result window.

\section{Relationship to the name service}

The node-local \texttt{ns} registry and replicated \texttt{dns} catalog are
separate services. DNS applies committed catalog state and resolves remote routes;
NS holds machine-local endpoint capabilities:

\begin{itemize}
    \item Cluster registrations are submitted through the Raft leader and applied
        to each DNS replica; local endpoint publication remains an NS operation.
    \item Service generations are replicated consistently. A restart on node 3
        cannot create a split-brain where node 1 sees generation 5 and node 2
        sees generation 3.
    \item Replicated policy metadata and tombstones propagate through the log;
        local delegation still depends on machine-local authority.
\end{itemize}

Kernel capability identifiers and endpoint objects are machine-local and are
not Raft state. For a local registration, lookup delegates a connection to the
local endpoint. For a remote registration, the local name-service instance
must return a local proxy connection whose transport reaches the owning node.
Raft replicates the metadata needed to choose that route, not the capability
object itself.

\section{Election timers as first-class completions}

In a traditional Raft implementation, election timeouts are implemented with
\texttt{sleep()} or a polling loop:

\begin{verbatim}
// Traditional: busy wait or sleep
while !timeout_elapsed() {
    sleep_ms(10); // or just spin
}
\end{verbatim}

In CharlotteOS, they use the completion system:

\begin{verbatim}
// CharlotteOS:
let timer_id = submit_timer(timeout_duration)?;
let result = wait_on_cq(timer_id).await?;
// Timer fired; start election.
\end{verbatim}

The timer is a kernel-managed operation. The shard parks on its CQ. When the
timer fires, the executor wakes, drains the CQ entry, and resumes the election
task. No polling loop, no thread dedicated to sleeping.

\section{Why consensus as a service matters for critical software}

\begin{enumerate}
    \item Consistency primitives can be exposed as services --- a replicated
        service implements \texttt{StateMachine} and uses the shared Raft,
        transport, timer, and storage interfaces.

    \item Capabilities, not IP addresses, identify peers --- the cluster
        configuration is a set of capabilities, not a list of \texttt{host:port}
        pairs. A network reconfiguration (new NIC, new subnet, move to RDMA)
        need not change the replicated peer identities.

    \item Timers are first-class and reliable --- an election timeout
        that is lost because the polling loop missed its window is a correctness
        bug (it can cause spurious leader changes or extended unavailability).
        Kernel-managed timers use the non-lossy CQ delivery path exercised by
        the local-election test.

    \item Failure isolation can extend to the consensus layer --- a Raft
        node is a userspace process. It can crash, restart, and recover without
        affecting the kernel or other Raft nodes. The supervisor handles the
        lifecycle; the Raft protocol handles the consensus.

    \item A distributed name service is implemented --- the
        \texttt{dns} service runs the \texttt{CatalogMachine} state machine
        on the Raft substrate; all other services that need cluster-wide
        consistency can use the same \texttt{StateMachine} pattern and shared
        consensus implementation.
\end{enumerate}

\endinput
